\documentclass[11pt,a4paper]{amsart}

\usepackage{cgeometry}

\addbibresource{ComplexGeometry.bib}
\defbibheading{bibliography}[\refname]{\section*{#1}}
\setlength{\bibitemsep}{0pt}
\renewcommand*{\bibfont}{\footnotesize}

\allowdisplaybreaks[2]
\emergencystretch=1em

%%%%%%%%%%%%%%%%%%%%%%%%%%%%%%%%%%%%%%%%%%%%%%%%%%%%%%%%%%%%%%%%%%%%%%%%%%%%%
%% Notation
%%%%%%%%%%%%%%%%%%%%%%%%%%%%%%%%%%%%%%%%%%%%%%%%%%%%%%%%%%%%%%%%%%%%%%%%%%%%%

\title{Real piecewise linearity of pluriharmonic functions on singular Berkovich spaces}
\author{}
\date{September 23, 2026}

\begin{document}

\begin{abstract}
Let $K$ be a complete nontrivially valued non-Archimedean field and let $X$ be a good, separated, strict, boundaryless $K$-analytic space. We prove that every classically pluriharmonic function on $X$ is locally real piecewise linear, including at non-quasi-smooth points. The descent argument uses closed-subspace rigidity for real unit-log formulas and regular alteration coverings of affinoid algebras.
\end{abstract}

\maketitle

\section{Introduction}\label{sec:introduction}

Let $X$ be a good strict analytic space over a complete nontrivially valued field $K$. A continuous function $h:X\to\mathbb R$ is real piecewise linear if it is G-locally a finite real linear combination of functions $\log|u|$, with $u$ an invertible analytic function. We write $\operatorname{PL}_{\mathbb R}$ for this sheaf. A function is classically pluriharmonic if it and its negative are classically plurisubharmonic in the curve-test sense, including the required tests after analytic scalar extension \cite[Definition~7.1]{GR25Classical}.

Real piecewise linearity is known when $X$ admits a strictly semistable formal model \cite[Proposition~7.8]{GR25Classical}. The quasi-smooth case follows by altered local uniformization and rigid \'{e}tale descent; we record the exact form used below in \Cref{prop:quasismooth-input}. The singular case asks whether the same conclusion holds without a prescribed finite flat or toroidal presentation. Such a presentation cannot simply be assumed for an arbitrary analytic germ.

\begin{maintheorem}\label{thm:singular-pluriharmonic}
Let $K$ be a complete nontrivially valued non-Archimedean field, and let $X$ be a good, separated, strict, boundaryless $K$-analytic space. Every classically pluriharmonic function $h:X\to\mathbb R$ is locally real piecewise linear. In particular this holds at every singular point of $X$.
\end{maintheorem}

The proof first establishes that a real PL germ on an analytic space is determined by its restriction to a closed analytic subspace through the point, and that every real PL germ on that subspace extends to the ambient space. These statements make it possible to descend PL functions through a proper modification once PL regularity is known off a closed subset and on that subset. Finite descent then follows by normalizing the base, treating finite normal extensions with a Galois and inseparable factorization, and using the conductor to recover the nonnormal base. A proper algebraic cover reduces to a finite cover over a dense open, and its exceptional subset is handled by dimension induction. Gabber's regular covering in the alteration topology and the domination theorem for that topology \cite[Expos\'{e}~VII, Th\'{e}or\`{e}me~1.1; Expos\'{e}~II, Th\'{e}or\`{e}me~3.2.1]{ILO14} supply the final regularizing geometry.

The proof uses relative analytification over a strict affinoid algebra. For a separated finite-type scheme $V$ over $\operatorname{Spec}(A)$, where $A$ is a strict $K$-affinoid algebra, write $V^{\mathrm{an}}$ for its relative Berkovich analytification over $\mathcal M(A)$. When discussing a PL germ at $x$, we use ordinary neighborhoods of $x$; its local formulas remain G-local on those neighborhoods. All closed analytic subspaces below carry their induced analytic structures.

\section{Closed analytic subspaces and real PL germs}\label{sec:closed-subspaces}

We record two affinoid facts used in the tube argument. A strict affinoid algebra admits a bounded Noether normalization, and the logarithmic image of finitely many analytic units on a strict affinoid space is a finite union of compact rational polytopes \cite[\S6.1]{BGR84}\cite[Theorem~A]{Ducros12Polytopes}. In the latter statement the polytope constants belong to $\mathbb Q\log|K^\times|$. We use rational domains with arbitrary positive real radii on nonstrict affinoids; their finite refinement property follows from the Gerritzen--Grauert theorem \cite{Temkin05GG}.

\begin{lemma}\label{lem:unit-norm-tube}
Let $Z=\mathcal M(A)$ be a strict $K$-affinoid space, let $T\hookrightarrow Z$ be closed analytic, and let $b\in A$. If $|b|\leq1$ on $T$, then $|b|\leq1$ on an ordinary open neighborhood of $T$ in $Z$. If $b$ is a unit and $|b|=1$ on $T$, the equality persists on such a neighborhood.
\end{lemma}

\begin{proof}
After replacing $T$ by its reduction, let $C$ be its reduced affinoid algebra. For each minimal prime $P$ of $C$, choose bounded Noether normalization $D=K\langle S_1,\ldots,S_d\rangle\hookrightarrow C/P$. The coordinates $S_i$ are images of functions on $Z$ of norm at most one. Put $F=\operatorname{Frac}(D)$ and $E=\operatorname{Frac}(C/P)$. The characteristic polynomial of multiplication by $b$ on the finite $F$-vector space $E$ has coefficients in $D$: they are integral over $D$ and lie in $F$, while $D$ is normal. At the Gauss point of $\mathcal M(D)$, its roots are values of $b$ on points of the finite fiber after a complete field extension. Their norms are at most one. The ultrametric inequality therefore bounds every coefficient by one at the Gauss point, hence in the Gauss norm of $D$.

Substitute the bounded lifts of the $S_i$ to obtain a monic polynomial $P_P(B)\in A[B]$ with coefficients of norm at most one and $P_P(b)=0$ on $\mathcal M(C/P)$. The product over the finitely many minimal primes is a monic polynomial $P(B)$ with bounded coefficients and $P(b)=0$ on $T$. On the ordinary open set $|P(b)|<1$, a value $|b|>1$ would make the leading term uniquely dominant, contradicting the inequality. This gives the first assertion. Apply it to $b$ and $b^{-1}$ for the second.
\end{proof}

\begin{lemma}\label{lem:real-log-tube}
Let $Z=\mathcal M(A)$ be any $K$-affinoid space and $T\hookrightarrow Z$ a closed analytic subspace. Suppose
\[
 f=\sum_{j=1}^{m}r_j\log|u_j|,\qquad r_j\in\mathbb R,\quad u_j\in A^\times,
\]
vanishes on $T$. Then $f$ vanishes on an ordinary open neighborhood of $T$ in $Z$.
\end{lemma}

\begin{proof}
First suppose $Z$ is strict. Fix $t\in T$ and let $P\subset\mathbb R^m$ be the logarithmic image of $(u_1,\ldots,u_m)$ on $T$. Write $w$ for the image of $t$ and $r=(r_1,\ldots,r_m)$. The hypothesis says $P\subset\{v:r\cdot v=0\}$. Near $w$, only the finitely many polytopes containing $w$ matter. Let $V$ be the sum of their rational direction spaces. Since all these polytopes contain $w$, their affine hulls lie in $w+V$, and $r$ annihilates $V$.

Choose integral row vectors $q_1,\ldots,q_s$ spanning the annihilator of $V$. On every polytope containing $w$, $q_i\cdot v$ has the constant value $c_i=q_i\cdot w\in\mathbb Q\log|K^\times|$. Choose $N_i>0$ and $a_i\in K^\times$ such that $\log|a_i|=N_ic_i$. On a strict affinoid neighborhood of $t$ contained in the inverse image of a small Euclidean ball around $w$, each ambient unit
\[
 b_i=a_i^{-1}\prod_j u_j^{N_iq_{ij}}
\]
has norm one on its intersection with $T$. By \Cref{lem:unit-norm-tube}, these equalities hold on an ordinary neighborhood of $t$. Since $r$ is a real linear combination of the $q_i$ and $r\cdot w=0$, they imply $f=0$ there. The union of these neighborhoods contains $T$.

For a general affinoid presentation with arbitrary radii, extend scalars to a complete analytic field $L/K$ whose value group contains those radii. The pullback $Z_L$ is strict, and the strict case gives a neighborhood of the full inverse image $T_L$ on which $f_L=0$. The projection $Z_L\to Z$ is a continuous surjection between compact Hausdorff spaces. The image of the compact complement of that neighborhood is closed and misses $T$; outside this image every lift lies in the vanishing neighborhood. Surjectivity gives the assertion on $Z$.
\end{proof}

\begin{proposition}\label{prop:closed-pl-germs}
Let $Z$ be a good separated $K$-analytic space and $T\hookrightarrow Z$ a closed analytic subspace. For every $t\in T$, restriction induces an isomorphism of real PL stalks
\[
 (\operatorname{PL}_{\mathbb R,Z})_t\longrightarrow(\operatorname{PL}_{\mathbb R,T})_t.
\]
\end{proposition}

\begin{proof}
For injectivity, represent a germ by a PL function vanishing on $T$ near $t$. A G-local affinoid unit-log presentation has, by the quasi-net property of analytic-domain coverings, finitely many compact affinoid charts covering an ordinary neighborhood of $t$. Shrink within each chart containing $t$ so that the formula vanishes on its closed subspace. By \Cref{lem:real-log-tube}, it vanishes on a relative ordinary neighborhood of $t$ in that chart. A chart missing $t$ and each compact complement of these relative neighborhoods can be avoided by one further ordinary shrinking. This gives zero on an ordinary neighborhood of $t$ in $Z$.

For surjectivity, represent the germ on a compact affinoid neighborhood $T_0$ of $t$ inside $T$. Choose a finite affinoid unit-log cover and refine it by rational domains $T_i$ of $T_0$. Lift the finitely many rational inequalities to an affinoid neighborhood $W$ of $T_0$ in $Z$. Their lifted rational domains $D_i$ cover a sufficiently small closed tube around $T_0$: the lifted Bezout relations generate the unit ideal on a neighborhood of $T_0$, and the nested closed tubes of the ideal of $T_0$ eventually lie in that neighborhood by compactness. Formation of rational domains commutes with quotient by the defining ideal, so $\mathcal O(D_i)\to\mathcal O(T_i)$ is surjective. Lift each analytic unit and its inverse in the formula on $T_i$. Their product equals one on $T_i$ and is a unit after shrinking the tube; the lift is therefore a unit on the corresponding $D_i$. The resulting lifted real unit-log functions agree on each overlap along $T_i\cap T_j$. By \Cref{lem:real-log-tube}, they agree on a neighborhood of that closed intersection. Finiteness and one more tube shrinking make the equalities hold throughout all overlaps. They glue to a PL function on the tube, whose ordinary interior contains $t$ and extends the given germ.

The same lifting construction works on nonstrict affinoid charts: use rational domains with arbitrary positive radii, whose coordinate rings still commute with closed-subspace quotients, and use the nonstrict case of \Cref{lem:real-log-tube} on their overlaps. All analytic units in this construction are lifted over $K$.
\end{proof}

\begin{proposition}\label{prop:proper-modification-descent}
Let $p:Y\to Z$ be a surjective morphism of good separated $K$-analytic spaces whose underlying continuous map is closed, let $T\hookrightarrow Z$ be closed analytic, and let $h:Z\to\mathbb R$ be continuous. If $h$ is locally real PL on $Z\setminus T$, $h|_T$ is locally real PL, and $h\circ p$ is locally real PL on $Y$, then $h$ is locally real PL on $Z$.
\end{proposition}

\begin{proof}
Fix $t\in T$. By surjectivity in \Cref{prop:closed-pl-germs}, extend the germ of $h|_T$ to a PL function $g$ on an ordinary neighborhood $U$ of $t$, with $g=h$ on $T\cap U$. The PL function $h\circ p-g\circ p$ vanishes on the closed analytic inverse image of $T\cap U$ in $p^{-1}(U)$. Injectivity in \Cref{prop:closed-pl-germs}, applied at each point of $p^{-1}(t)$, makes the difference zero on an ordinary neighborhood $V$ of that entire fiber in $p^{-1}(U)$. The restricted map $p^{-1}(U)\to U$ is closed, so $U'=U\setminus p(p^{-1}(U)\setminus V)$ is an ordinary neighborhood of $t$ and $p^{-1}(U')\subset V$. Surjectivity gives $h=g$ on $U'$. Points outside $T$ were covered by the first hypothesis.
\end{proof}

\section{Finite descent}\label{sec:finite-descent}

The first two lemmas replace a fiber-transitivity argument for an arbitrary finite quotient. Their proofs keep all subdivisions rational even though the coefficients of a PL function may be irrational.

\begin{lemma}\label{lem:bounded-tropical-presentation}
Let $Z=\mathcal M(A)$ be strict affinoid and let $f\in\operatorname{PL}_{\mathbb R}(Z)$. There are $a_1,\ldots,a_m\in A$, a positive $\delta\in|K^\times|$, finitely many closed polyhedra $R_i\subset\mathbb R^m$ cut out by rational linear inequalities with real constants, and real affine functions $L_i$ such that, with
\[
 \tau(z)=\bigl(\log\max(|a_1(z)|,\delta),\ldots,\log\max(|a_m(z)|,\delta)\bigr),
\]
the sets $\tau^{-1}(R_i)$ cover $Z$ and $f=L_i\circ\tau$ on each of them. The $L_i$ agree on overlaps in $\tau(Z)$.
\end{lemma}

\begin{proof}
Take a finite rational-domain refinement $D_i$ of a G-local affinoid unit-log cover. Write the inequalities of $D_i$ as $|p_{ij}|\leq r_{ij}|g_i|$. The functions $g_i,p_{ij}$ generate the unit ideal, and $|g_i|$ has a positive minimum on $D_i$. An analytic unit $u$ on $D_i$ can be uniformly approximated by $a/g_i^N$, with $a\in A$, to an error smaller than $\min_{D_i}|u|$. The ultrametric inequality gives $|u|=|a/g_i^N|$ on $D_i$. Consequently the unit-log formula for $f$ is a real linear combination of logarithms of finitely many ambient functions nonvanishing on $D_i$.

Collect these functions, the $g_i$, and the $p_{ij}$ into one finite list $a_l$. For each $i$, the functions $g_i,p_{ij}$ generate the unit ideal, so their maximum norm has a positive lower bound $c_i$ on $Z$. Choose $\delta\in|K^\times|$ below $c_i/\max(1,r_{ij})$ for every $i,j$, below every positive lower bound used in the formulas, and below $r_{ij}\min_{D_i}|g_i|$. If $|g_i|<\delta$, some $|p_{ij}|$ exceeds the corresponding truncated right side; hence the truncated inequalities cannot admit a point outside $D_i$ through a common small denominator. Thus the inequalities defining each $D_i$ are equivalent to rational affine inequalities in the coordinates of $\tau$, and truncation at $\delta$ does not alter a formula on its domain. These give $R_i$ and $L_i$. Two points with the same $\tau$-vector satisfy the same domain inequalities, so the formulas agree on overlaps in $\tau(Z)$.
\end{proof}

\begin{lemma}\label{lem:support-vector-selection}
For integers $m,d\geq1$ there are finitely many $q^j\in\mathbb Z_{>0}^m$ and an integer $N\geq1$ with the following property. Given $v_1,\ldots,v_d\in\mathbb R^m$, put $H(a)=\max_s a\cdot v_s$. For some $j$,
\begin{equation}\label{eq:support-vector-equality}
 H(Nq^j+e_l)+H(Nq^j-e_l)=2NH(q^j)\qquad(1\leq l\leq m),
\end{equation}
and whenever these equalities hold, the vector with $l$-th coordinate $H(Nq^j+e_l)-NH(q^j)$ is one of the $v_s$. All vectors $Nq^j\pm e_l$ can be chosen nonnegative integral.
\end{lemma}

\begin{proof}
The case $d=1$ is immediate. Otherwise put $M=d(d-1)/2$. Choose more than $M(m-1)$ distinct positive integer vectors on the moment curve $(1,t,\ldots,t^{m-1})$. A nonzero linear functional vanishes on at most $m-1$ of them. Compactness of the product of $M$ unit spheres gives a constant $c>0$ such that, for any list of at most $M$ nonzero differences $v_s-v_t$, some chosen $q^j$ satisfies $|q^j\cdot(v_s-v_t)|\geq c\|v_s-v_t\|$ for all pairs with distinct vectors. Choose $N$ with $Nc\geq1$ and enlarge it so that $Nq^j-e_l$ is nonnegative. The maximizer of $q^j\cdot v_s$, unique up to repeated identical vectors, also maximizes each $(Nq^j\pm e_l)\cdot v_s$, proving \eqref{eq:support-vector-equality}. Conversely, if the equalities hold, take any maximizer $v_*$ of $q^j\cdot v_s$. Each term on the left is at least $NH(q^j)\pm(v_*)_l$; equality of the sums makes both bounds equalities. The recovered vector is $v_*$.
\end{proof}

\begin{proposition}\label{prop:finite-group-descent}
Let $A\to B$ be finite between strict $K$-affinoid algebras. Suppose a finite group $G$ acts on $B$ by $A$-algebra automorphisms, $B^G=A$, and $\mathcal M(B)\to\mathcal M(A)$ is surjective. If a continuous function $h$ on $\mathcal M(A)$ becomes real PL on $\mathcal M(B)$, then $h$ is real PL.
\end{proposition}

\begin{proof}
Apply \Cref{lem:bounded-tropical-presentation} to the pullback $f$ of $h$, with ambient functions $b_1,\ldots,b_m\in B$ and $\delta=|a|$ for $a\in K^\times$. For a nonnegative integral vector $q$, and each $0\leq\nu_l\leq q_l$, set $c_\nu=a^{\sum_l(q_l-\nu_l)}\prod_l b_l^{\nu_l}$. If $|G|=d$, write
\[
 \prod_{g\in G}(T-g(c_\nu))=T^d+\sum_{k=1}^d e_{\nu,k}T^{d-k}.
\]
Every coefficient $e_{\nu,k}$ lies in $B^G=A$. Define on $\mathcal M(A)$
\[
 M_q(x)=\log\max_{\nu,k}|e_{\nu,k}(x)|^{1/k}.
\]
This is finite and bounded below by $\delta^{\sum q_l}$ inside the logarithm, using $\nu=0$. The non-Archimedean root-radius identity, evaluated at any $y$ over $x$, yields
\begin{equation}\label{eq:invariant-support-function}
 M_q(x)=\max_{g\in G}\sum_{l=1}^m q_l\log\max(|g(b_l)(y)|,\delta).
\end{equation}
Thus the right side is independent of the chosen lift without any assertion about transitivity of $G$ on Berkovich fibers.

Each $M_q$ has a finite rational cover on which it is $(1/k)\log|e_{\nu,k}|$ for a winning pair $(\nu,k)$. Indeed, compare candidate coefficients after clearing the integer roots, and impose the positive lower bound in the same rational domain; the winning coefficient is then a unit. Fix the finite collection of vectors $q^j,Nq^j\pm e_l$ in \Cref{lem:support-vector-selection}, and take a common such rational refinement for their $M$-functions. For a chosen lift $y$, the vectors $v_g=(\log\max(|g(b_l)(y)|,\delta))_l$ satisfy $M_q(x)=\max_g q\cdot v_g$ by \eqref{eq:invariant-support-function}. The equalities \eqref{eq:support-vector-equality} define finitely many closed subsets covering $\mathcal M(A)$. On the subset for $j$, the coordinates
\[
 w_l=M_{Nq^j+e_l}-NM_{q^j}
\]
form one of the $v_g$. Refine once more by the rational inequalities defining the polyhedra $R_i$ of \Cref{lem:bounded-tropical-presentation}, applied to $w$. On each resulting rational domain, $h=L_i(w)$ is a real linear combination of logarithms of the winning invariant coefficients, plus a real constant. Such a constant is real PL because the valuation is nontrivial. These domains cover the base and prove the assertion.
\end{proof}

\begin{lemma}\label{lem:frobenius-descent}
Suppose $A\hookrightarrow B$ is a finite radicial surjective map of strict affinoid algebras in characteristic $p>0$ and $B^q\subset A$ for a fixed power $q=p^e$. If a continuous function $h$ on $\mathcal M(A)$ pulls back to real PL on $\mathcal M(B)$, then $h$ is real PL.
\end{lemma}

\begin{proof}
The finite radicial map is a homeomorphism on analytic spectra. Refine the upstairs PL presentation by finitely many rational domains. The $q$-th powers of their defining functions and of the analytic units in each formula belong to the corresponding descended rational-domain algebra. Taking $q$-th powers of the inequalities gives a rational cover downstairs, and $q$-th powers of a unit and its inverse remain inverse there. Dividing every real coefficient by $q$ gives the same function on the descended domain. The homeomorphism makes these domains cover $\mathcal M(A)$.
\end{proof}

\begin{proposition}\label{prop:normal-finite-descent}
Let $A$ be a normal integral strict $K$-affinoid algebra and $B$ a finite strict $K$-affinoid $A$-algebra for which $\mathcal M(B)\to\mathcal M(A)$ is surjective. If a continuous function on $\mathcal M(A)$ becomes real PL on $\mathcal M(B)$, then it is real PL on $\mathcal M(A)$.
\end{proposition}

\begin{proof}
Replace $B$ by a minimal-prime quotient dominating $A$. Such a quotient exists: otherwise a product of nonzero elements in the finitely many contractions would be nonzero in the domain $A$ but nilpotent in $B$. This replacement preserves finite surjectivity and the pulled-back PL function. Put $F=\operatorname{Frac}(A)$ and $E=\operatorname{Frac}(B)$. Choose a finite normal extension $L/F$ containing $E$, let $L_s/F$ be its maximal separable subextension, and let $C$ and $D$ be the integral closures of $A$ in $L$ and $L_s$. Excellence of strict affinoid algebras makes $C$ and $D$ finite over $A$ \cite[proof of Theorem~1.3.1]{Berkovich15Finiteness}. The pullback to $\mathcal M(C)$ is PL. In positive characteristic a fixed $p$-power sends $C$ into $D$, so \Cref{lem:frobenius-descent} gives PL on $\mathcal M(D)$. In characteristic zero this step is empty. The finite Galois group $G=\operatorname{Gal}(L_s/F)$ acts on $D$, and $D^G=A$ by normality of $A$. Apply \Cref{prop:finite-group-descent} to $D/A$.
\end{proof}

\begin{theorem}\label{thm:finite-descent}
Let $A\to B$ be finite between strict $K$-affinoid algebras, with $\mathcal M(B)\to\mathcal M(A)$ surjective. If a continuous real function on $\mathcal M(A)$ pulls back to real PL on $\mathcal M(B)$, then it is real PL on $\mathcal M(A)$.
\end{theorem}

\begin{proof}
Nilpotents can be discarded using the closed-immersion stalk isomorphism in \Cref{prop:closed-pl-germs}. We induct on $d=\dim A$ for reduced $A$. Let $C=\prod_j C_j$ be its normalization, finite over $A$ by affinoid excellence. Each $C_j$ is normal integral, and $B\otimes_A C_j$ is finite surjective over $C_j$. By \Cref{prop:normal-finite-descent}, $h$ pulls back to a PL function on $\mathcal M(C)$.

Put $I=\operatorname{Ann}_A(C/A)$ and $T=\mathcal M(A/I)$. The finite module $C/A$ vanishes at each minimal prime of $A$, so $I$ avoids every minimal prime and contains a nonzerodivisor by prime avoidance. Hence $\dim(A/I)<d$ when $T$ is nonempty. Normalization is an isomorphism off $T$, making $h$ PL there. The restriction $h|_T$ becomes PL after the finite surjective base change $B\otimes_A A/I$; induction makes $h|_T$ PL. Apply \Cref{prop:proper-modification-descent} to the finite normalization map $\mathcal M(C)\to\mathcal M(A)$. In dimension zero the reduced algebra is a product of fields and equals its normalization, which starts the induction.
\end{proof}

\section{Proper algebraic descent}\label{sec:proper-descent}

For a strict affinoid algebra $A$, relative analytification of a separated finite-type $A$-scheme commutes with base change to a strict affinoid algebra, with open and closed immersions, and with finite morphisms. A proper $A$-scheme has compact analytification over $\mathcal M(A)$ by the properness comparison \cite[Proposition~2.6.9]{Berkovich93Etale}; this compactness is used explicitly in \cite[proof of Theorem~1.3.1]{Berkovich15Finiteness}. A proper surjective scheme morphism gives a surjective analytic map: over $x\in\mathcal M(A)$ its nonempty finite-type fiber over $\mathscr H(x)$ has a closed point with finite residue extension, and the unique extension of the absolute value to that field supplies an analytic lift of $x$. The compact analytic source and Hausdorff target make this map closed. These facts also apply to proper morphisms between separated finite-type $A$-schemes after relative analytification.

\begin{lemma}\label{lem:generic-multisection}
Let $A$ be a Noetherian domain and let $V\to\operatorname{Spec}(A)$ be proper and surjective. There is an integral closed subscheme $W\subset V$ which is proper and surjective over $\operatorname{Spec}(A)$ and a nonzero $c\in A$ such that $W_{A[1/c]}\to\operatorname{Spec}(A[1/c])$ is finite surjective.
\end{lemma}

\begin{proof}
Write $F=\operatorname{Frac}(A)$. The nonempty finite-type $F$-scheme $V_F$ has a closed point $w$ with finite residue extension $E/F$. Let $W$ be its reduced closure in $V$. It is integral and proper over $A$, and its generic fiber is $\operatorname{Spec}(E)$. Its proper image is closed and contains the generic point of $\operatorname{Spec}(A)$, so it is surjective.

Choose an affine open $W_0=\operatorname{Spec}(R)$ containing $w$. The proper image of the closed complement $W\setminus W_0$ misses the generic point. After inverting a nonzero $b\in A$, the whole $W_{A[1/b]}$ equals $(W_0)_{A[1/b]}$. The finite-type $A$-algebra $R$ is a domain with $R\otimes_A F=E$. Choose finite algebra generators $r_i$ of $R$. Each satisfies a monic equation over $F$; after inverting one further nonzero $a\in A$, all coefficients lie in $A[1/a]$. Since $R$ embeds into $E$, these equations hold in $R[1/a]$, making it a finite $A[1/a]$-module. Set $c=ab$. Surjectivity persists after this base change.
\end{proof}

\begin{proposition}\label{prop:proper-algebraic-descent}
Let $A$ be a strict $K$-affinoid algebra and let $V\to\operatorname{Spec}(A)$ be proper and surjective. If $h:\mathcal M(A)\to\mathbb R$ is continuous and its pullback to $V^{\mathrm{an}}$ is locally real PL, then $h$ is real PL.
\end{proposition}

\begin{proof}
We induct on $d=\dim A$. Removing nilpotents does not change the question by \Cref{prop:closed-pl-germs}. For a reduced $A$, the finite surjection $\coprod_P\operatorname{Spec}(A/P)\to\operatorname{Spec}(A)$ over its minimal primes lets us first prove the assertion for each integral quotient: properness and surjectivity survive that base change, and \Cref{thm:finite-descent} then returns PL to $\mathcal M(A)$. Thus assume $A$ integral.

Choose $W$ and $c$ by \Cref{lem:generic-multisection}. The function pulls back to a PL function on $W^{\mathrm{an}}$. Over $U=\{x:|c(x)|>0\}$ the map $W^{\mathrm{an}}\to\mathcal M(A)$ is finite surjective, as can be checked after restricting to strict affinoid domains in $U$. By \Cref{thm:finite-descent}, $h$ is locally PL on $U$. If $S=\mathcal M(A/(c))$ is nonempty, it has dimension less than $d$. The base change of $V$ to $A/(c)$ remains proper surjective and the pullback of $h|_S$ remains PL, so induction gives PL for $h|_S$. The map $W^{\mathrm{an}}\to\mathcal M(A)$ is closed and surjective by the analytification facts above. Apply \Cref{prop:proper-modification-descent} to $S$ to conclude. In dimension zero $c$ is a unit and the finite step alone starts the induction.
\end{proof}

\section{Regular covers and the quasi-smooth case}\label{sec:regular-covers}

The next proposition records the quasi-smooth input with its descent mechanism. Altered local uniformization and the semistable PL theorem both apply on compact strict spaces, including their analytic boundaries. This is needed when they are applied to an affinoid neighborhood.

\begin{proposition}\label{prop:quasismooth-input}
Let $K$ be algebraically closed, complete, and nontrivially valued. If $Y$ is a good, separated, strict, quasi-smooth $K$-analytic space, every classically pluriharmonic function on $Y$ is locally real PL.
\end{proposition}

\begin{proof}
The assertion is local on $Y$, so take a compact strict affinoid neighborhood $Z$ of a point. The associated rigid space is smooth, quasi-compact, and quasi-separated. Temkin's altered local uniformization gives a rigid \'{e}tale covering by affinoids with strictly semistable affine formal models after a finite field extension \cite[Corollary~3.3.2]{Temkin17}. Since $K$ is algebraically closed, that extension is $K$. The pullback of the given function to each member is classically pluriharmonic by functoriality \cite[Proposition~7.2(7)]{GR25Classical}. It is real PL by \cite[Proposition~7.8]{GR25Classical}.

The descent step uses the rigid \'{e}tale sheaf criterion in \cite[Proposition~3.2.2]{dJvdP96}: a presheaf is a sheaf for rigid \'{e}tale coverings if it is a sheaf for admissible rigid-domain coverings and satisfies the equalizer condition for finite \'{e}tale surjections of affinoids. The real PL presheaf has the first property by its G-local definition. For the second, compatible PL functions on a finite \'{e}tale cover determine a continuous function downstairs. Indeed, two points in one fiber admit a common point of the analytic fiber product, obtained from a residue field of the tensor product of their finite completed residue fields over the base residue field. Compatibility therefore makes their values equal. A finite surjection is a closed quotient map, so the induced function downstairs is continuous. It is PL by \Cref{thm:finite-descent}, applied on strict affinoid neighborhoods. Thus the semistable pullbacks descend to a real PL function on $Z$, equal to the original function by surjectivity.
\end{proof}

\begin{lemma}\label{lem:regular-analytification}
Let $K$ be algebraically closed, complete, and nontrivially valued, let $A$ be a strict $K$-affinoid algebra, and let $R$ be a regular separated finite-type scheme over $\operatorname{Spec}(A)$. Then $R^{\mathrm{an}}$ is a good, separated, strict, quasi-smooth $K$-analytic space.
\end{lemma}

\begin{proof}
The first three properties follow from relative analytification. Quasi-smoothness can be tested on strict affinoid charts and is detected there by the analytic Jacobian criterion. Cover $R$ by affine charts $\operatorname{Spec}(A[T_1,\ldots,T_n]/J)$, and choose a strict presentation $A=K\langle X_1,\ldots,X_m\rangle/I$. At a rigid point of an analytified chart, algebraic closedness makes both residue fields equal to $K$. Choose a polydisc chart with the point in its relative interior. The completed local rings of the algebraic and analytic charts at that point are both the quotient of $K[[\xi_1,\ldots,\xi_m,\tau_1,\ldots,\tau_n]]$ by the completed equations from $I$ and $J$. The algebraic local ring is regular, so its completion is regular. The Jacobian criterion at the $K$-rational point therefore makes the analytic chart quasi-smooth there, in any characteristic. On a strict affinoid chart the non-quasi-smooth locus is closed analytic and, if nonempty, has a rigid point by the affinoid Nullstellensatz. Hence it is empty.
\end{proof}

\begin{lemma}\label{lem:regular-proper-cover}
Let $A$ be an integral strict $K$-affinoid algebra. There exist an integral $A$-scheme $V$, proper and surjective over $\operatorname{Spec}(A)$, a finite Zariski open cover $V=\bigcup_j V_j$, and regular integral separated finite-type $A$-schemes $R_j$ with $A$-morphisms $V_j\to R_j$.
\end{lemma}

\begin{proof}
Strict affinoid algebras are excellent \cite[proof of Theorem~1.3.1]{Berkovich15Finiteness}, so $S=\operatorname{Spec}(A)$ is noetherian and quasi-excellent. Choose a prime $\ell$ invertible on $S$. Gabber's weak local uniformization supplies a finite $\ell'$-alteration covering family $R_i\to S$, with each $R_i$ regular and connected \cite[Expos\'{e}~VII, Th\'{e}or\`{e}me~1.1, with $Z=\varnothing$]{ILO14}. Forgetting the prime-to-$\ell$ degree condition makes it a covering family in the ordinary alteration topology. A noetherian regular connected scheme is integral. Refine each $R_i$ by a finite affine open cover. This remains a covering family in the same topology and makes its regular members separated and of finite type over $S$. For irreducible noetherian $S$, the domination theorem for that topology gives an integral proper surjection $V\to S$ and a Zariski open cover $V=\bigcup_jV_j$ such that each composite $V_j\to S$ factors through one of these regular affine members \cite[Expos\'{e}~II, Th\'{e}or\`{e}me~3.2.1, case $\ell=1$]{ILO14}. Rename the corresponding regular schemes $R_j$.
\end{proof}

\begin{proposition}\label{prop:algebraically-closed-case}
Let $K$ be algebraically closed, complete, and nontrivially valued, let $A$ be a strict $K$-affinoid algebra, and let $h$ be a continuous classically pluriharmonic function on $\mathcal M(A)$. Then $h$ is real PL.
\end{proposition}

\begin{proof}
First suppose $A$ integral. Choose $V$, $V_j$, and $R_j$ by \Cref{lem:regular-proper-cover}. Classical pluriharmonicity is preserved by analytic pullback, since every curve test on a source composes with the map to $\mathcal M(A)$, also after analytic field extension. By \Cref{lem:regular-analytification,prop:quasismooth-input}, the pullback of $h$ to each $R_j^{\mathrm{an}}$ is locally real PL. Its pullback to $V_j^{\mathrm{an}}$ is therefore locally real PL. The sets $V_j^{\mathrm{an}}$ are ordinary open in $V^{\mathrm{an}}$ and cover it; hence the pullback to $V^{\mathrm{an}}$ is locally real PL. \Cref{prop:proper-algebraic-descent} gives real PL on $\mathcal M(A)$.

For arbitrary $A$, pass to its reduction and then to the finite product of its minimal-prime quotients. On each integral quotient the pulled-back function is classically pluriharmonic and hence real PL by the preceding paragraph. \Cref{thm:finite-descent} returns real PL to $\mathcal M(A_{\mathrm{red}})$, and \Cref{prop:closed-pl-germs} identifies its PL germs with those of $\mathcal M(A)$.
\end{proof}

\section{Scalar descent and proof of the theorem}\label{sec:scalar-descent}

\begin{proposition}\label{prop:completed-algebraic-descent}
Let $L$ be the completion of an algebraic closure of $K$, let $Z=\mathcal M(A)$ be strict $K$-affinoid, and let $h:Z\to\mathbb R$ be continuous. If the pullback of $h$ to the whole space $Z_L$ is real PL, then $h$ is real PL on $Z$.
\end{proposition}

\begin{proof}
By compactness, take a finite rational-domain cover of $Z_L$ on which the pullback has unit-log formulas. In a rational presentation $|f_j|\leq r_j|g|$, the functions $f_j,g$ generate the unit ideal, so their maximum norm has a positive lower bound $c$ on $Z_L$. Uniform approximation of all these functions by elements defined over finite intermediate extensions of $K$ changes neither the rational set nor its analytic-domain structure when every error is smaller than $c$ and the corresponding $r_jc$ bounds: on either rational set the denominator is bounded below, and ultrametricity preserves each defining inequality. Since the algebraic elements of $L$ are dense, one finite extension $F/K$ inside $L$ descends the entire finite rational cover. The unit functions on these domains can likewise be approximated by functions over a finite intermediate field, with errors smaller than their positive lower norms. Ultrametricity preserves their absolute values; the approximants are units because every downstairs point has a lift to $L$. After enlarging $F$ finitely many times, the formulas descend to a PL function on $Z_F$. On overlaps they agree after scalar extension to $L$, hence already on $Z_F$. Their pullback to $Z_L$ is the pullback of $h$, so surjectivity gives the function $h_F$ itself. The finite surjection $Z_F\to Z$ now descends PL by \Cref{thm:finite-descent}.
\end{proof}

\begin{proof}[Proof of \Cref{thm:singular-pluriharmonic}]
Both $h$ and $-h$ are upper semicontinuous by definition, so $h$ is continuous. Let $x\in X$ and choose a strict affinoid ordinary neighborhood $Z=\mathcal M(A)$ of $x$ in $X$. Let $L$ be the completion of an algebraic closure of $K$; it is algebraically closed. The pullback $h_L$ is classically pluriharmonic on $Z_L$ by scalar-extension invariance and functoriality \cite[Proposition~7.2(6)--(7)]{GR25Classical}. By \Cref{prop:algebraically-closed-case}, $h_L$ is real PL on the entire $Z_L$. \Cref{prop:completed-algebraic-descent} makes $h$ real PL on $Z$, and restriction to the ordinary interior of $Z$ proves the assertion near $x$.
\end{proof}

\printbibliography

\end{document}
